\documentclass[11pt,a4paper]{article}
\usepackage[margin=2.2cm,top=2.5cm,bottom=2.5cm]{geometry}
\usepackage{amsmath,amssymb,amsfonts}
\usepackage{booktabs}
\usepackage{tabularx}
\usepackage{listings}
\usepackage{xcolor}
\usepackage{hyperref}
\usepackage{graphicx}
\usepackage{parskip}
\usepackage{microtype}
\usepackage{array}
\usepackage{cite}
\usepackage{caption}
\usepackage{subcaption}
\usepackage{enumitem}
\usepackage{fancyhdr}
\usepackage{titlesec}
\usepackage[most]{tcolorbox}

% ── Color Palette ─────────────────────────────────────────────────────────────
\definecolor{primary}{RGB}{24, 76, 120}     % Deep academic navy
\definecolor{secondary}{RGB}{41, 128, 185}  % Bright blue accent
\definecolor{darkslate}{RGB}{44, 62, 80}    % Charcoal body
\definecolor{codebg}{RGB}{248, 249, 250}    % Clean subtle light grey
\definecolor{codeborder}{RGB}{220, 224, 230}
\definecolor{boxbg}{RGB}{243, 247, 252}
\definecolor{boxborder}{RGB}{180, 205, 235}
\definecolor{kwcolor}{RGB}{0, 90, 160}
\definecolor{comcolor}{RGB}{110, 120, 135}
\definecolor{strcolor}{RGB}{180, 40, 40}

% ── Hyperref Setup ────────────────────────────────────────────────────────────
\hypersetup{
    colorlinks=true,
    linkcolor=secondary,
    citecolor=primary,
    urlcolor=secondary,
    pdfauthor={Sanan Akther, Divyansh Sharma, Joseph Maher Labib Estawro},
    pdftitle={High-Performance Iterative Cone-Beam CT Volume Reconstruction},
    pdfkeywords={CT Reconstruction, MLEM, OSEM, OpenCL, GPU Acceleration, HPC}
}

% ── Header & Footer Setup ─────────────────────────────────────────────────────
\pagestyle{fancy}
\fancyhf{}
\fancyhead[L]{\small\color{gray}\nouppercase{\leftmark}}
\fancyhead[R]{\small\color{gray}GPULAB — Final Course Report}
\fancyfoot[C]{\small\thepage}
\renewcommand{\headrulewidth}{0.4pt}
\renewcommand{\headrule}{\hbox to\headwidth{\color{codeborder}\leaders\hrule height \headrulewidth\hfill}}

% ── Section Heading Styling ───────────────────────────────────────────────────
\titleformat{\section}
  {\Large\bfseries\color{primary}}
  {\thesection}{1em}{}
  [{\color{secondary}\titlerule[0.8pt]}]

\titleformat{\subsection}
  {\large\bfseries\color{darkslate}}
  {\thesubsection}{0.8em}{}

\titleformat{\subsubsection}
  {\normalsize\bfseries\color{darkslate}}
  {\thesubsubsection}{0.6em}{}

% ── Listings Setup ────────────────────────────────────────────────────────────
\lstset{
  basicstyle=\ttfamily\scriptsize,
  keywordstyle=\color{kwcolor}\bfseries,
  commentstyle=\color{comcolor}\itshape,
  stringstyle=\color{strcolor},
  numberstyle=\tiny\color{gray!80},
  numbers=left,
  stepnumber=1,
  numbersep=6pt,
  frame=single,
  frameround=tttt,
  rulecolor=\color{codeborder},
  backgroundcolor=\color{codebg},
  breaklines=true,
  breakatwhitespace=true,
  tabsize=2,
  language=C,
  captionpos=b,
  showstringspaces=false,
  xleftmargin=12pt,
  framexleftmargin=8pt
}

% ── Table Column Types (Auto-wrapping & Centering) ────────────────────────────
\newcolumntype{Y}{>{\centering\arraybackslash}X}
\newcolumntype{L}{>{\raggedright\arraybackslash}X}
\newcolumntype{C}[1]{>{\centering\arraybackslash}p{#1}}
\newcolumntype{R}[1]{>{\raggedleft\arraybackslash}p{#1}}

% ── Callout Box Environment ───────────────────────────────────────────────────
\newtcolorbox{infobox}[1][]{
  colback=boxbg,
  colframe=boxborder,
  arc=3mm,
  boxrule=0.8pt,
  left=3mm, right=3mm, top=2mm, bottom=2mm,
  title=\textbf{#1},
  coltitle=primary,
  fonttitle=\small\bfseries,
  attach boxed title to top left={yshift=-2mm, xshift=3mm},
  boxed title style={colback=boxbg, colframe=boxborder, boxrule=0.5pt, arc=1mm}
}

% ── Title & Author Block ──────────────────────────────────────────────────────
\title{
  \vspace{-1.5cm}
  {\Huge\textbf{\color{primary}High-Performance Iterative Cone-Beam CT}}\\[0.3em]
  {\Huge\textbf{\color{primary}Volume Reconstruction}}\\[0.5em]
  {\large\textbf{\color{darkslate}GPULAB: High Performance Computing with Graphic Cards}}\\[0.2em]
  {\normalsize\color{gray}Final Project Report — Topic 1-2: CT Volume Reconstruction}
  \vspace{-0.3cm}
}

\author{
  \textbf{Sanan Akther} \quad (Matriculation No. 3704154)\\[0.25em]
  \textbf{Divyansh Sharma} \quad (Matriculation No. 3758980)\\[0.25em]
  \textbf{Joseph Maher Labib Estawro} \quad (Matriculation No. 3699049)\\[0.8em]
  \normalsize Institute of Computer Architecture and Computer Engineering (ITI)\\
  \normalsize University of Stuttgart, Germany
}
\date{September 2026}

\begin{document}
\maketitle
\thispagestyle{plain}

\begin{abstract}
\noindent
Iterative statistical reconstruction for cone-beam Computed Tomography (CBCT) provides superior image fidelity over analytic filtered backprojection (FDK) under noisy or sparse-view acquisitions, but its computational intensity poses a severe latency bottleneck. In this project, we develop, optimize, and evaluate a high-performance heterogeneous CPU and GPU reconstruction framework based on the Maximum Likelihood Expectation Maximization (MLEM) and Ordered Subsets EM (OSEM) algorithms. Starting from an unaccelerated Python reference, we implement a multi-threaded C/OpenMP CPU solver and three distinct OpenCL GPU pipelines: a naive global buffer mode (\texttt{gpu-buf}), a hardware texture-sampler-accelerated image mode (\texttt{gpu-img}), and an optimized pipeline (\texttt{gpu-opt}) featuring precomputed trigonometric lookup tables, cooperative local memory caching, and fused elementwise kernels.

Through rigorous numerical audits, we identify and resolve critical correctness discrepancies, including fractional index truncation bugs and whole-cell interpolation boundary conditions, bringing CPU--GPU agreement to the $\text{float32}$ precision limit ($\text{MSE} \approx 10^{-10}\text{--}10^{-7}$). On the CPU, an 8-way ray-tiling memory reordering (\texttt{FP\_TILE=32}) accelerates forward projection by $>2\times$ by eliminating cache scatter across multi-megabyte strides. On the GPU, hardware texture sampling combined with Phase~B kernel fusions (\texttt{divide\_preprocess\_img} and \texttt{vol\_update\_img}) eliminates intermediate buffer roundtrips, while systematic work-group sweeps ($\{8,32,1\}$ for projection ray-marching and $\{4,4,16\}$ for coalesced backprojection writes) deliver up to $\approx 30\times$ speedup over the 12-thread CPU baseline. Furthermore, we formulate and integrate Ordered Subsets EM (OSEM) with golden-ratio-derived coprime subset ordering, demonstrating continuous objective optimization that overcomes the plateauing behavior of plain MLEM within fixed wall-clock time budgets. All implementations are validated across full 100-epoch horizons on both $256^3$ and $512^3$ volumes across AMD GCN (Hawaii PRO) and NVIDIA Kepler (GTX 680) architectures.
\end{abstract}

\newpage
\begingroup
\hypersetup{linkcolor=darkslate}
\setlength{\parskip}{0.5pt}
\setcounter{tocdepth}{2}
\tableofcontents
\endgroup
\newpage

% ─────────────────────────────────────────────────────────────────────────────
\section{Introduction and Problem Formulation}
\label{sec:intro}

Cone-beam Computed Tomography (CBCT) is a fundamental imaging modality across clinical radiotherapy planning, image-guided surgery, dental diagnostics, and non-destructive industrial testing. In CBCT, an X-ray point source and a two-dimensional flat-panel matrix detector rotate synchronously around an object along a circular gantry trajectory, capturing a sequence of cone-beam transmission radiographs at discrete projection angles $\theta$.

\subsection{Physical Geometry}
The acquisition setup is parameterized by two fundamental geometric distances:
\begin{itemize}[noitemsep,topsep=2pt]
    \item \textbf{Source-to-Object Distance ($\text{SOD}$)}: The radial distance from the focal spot of the X-ray point source to the center of rotation (isocenter).
    \item \textbf{Source-to-Detector Distance ($\text{SDD}$)}: The total distance from the X-ray point source to the flat detector plane.
\end{itemize}
The 3D volume attenuation field is discretized on a Cartesian voxel grid $v(x,y,z)$ of dimensions $N_{xz} \times N_{xz} \times N_y$ with isotropic voxel spacing $\delta_v$. The flat-panel detector has dimensions $W \times H$ pixels with physical pixel spacing $\delta_p$, indexed by $(u, v)$ coordinates centered at the detector optical axis.

\subsection{Mathematical Model of Continuous and Discrete Operators}
For a projection angle $\theta$, the continuous Radon line integral through attenuation field $v(\mathbf{r})$ along a ray $\mathcal{L}(u, v, \theta)$ parameterized by path length $t$ is:
\begin{equation}
    p(u, v, \theta) = \int_{\mathcal{L}(u, v, \theta)} v(\mathbf{r}(t)) \, \mathrm{d}t
\end{equation}

\subsubsection{Forward Projection Operator ($\mathcal{F}$)}
In the discrete computational model, the forward projection operator $\mathcal{F}: \mathbb{R}^{N_{xz} \times N_{xz} \times N_y} \to \mathbb{R}^{N_\theta \times H \times W}$ computes the line integral for each detector pixel via numerical quadrature. Using incremental ray-marching with $N_s$ equidistant samples along the ray between entry $t_{\text{near}}$ and exit $t_{\text{far}}$:
\begin{equation}
    b(u, v, \theta) = \mathcal{F}(v)(u, v, \theta) \approx \sum_{s=0}^{N_s-1} v\big(\mathbf{r}(t_{\text{near}} + s \cdot \Delta t)\big) \cdot \Delta t \cdot \|\mathbf{d}_{\text{ray}}\|_2
    \label{eq:fp_discrete}
\end{equation}
where $\mathbf{r}(t) = \mathbf{S}(\theta) + t \cdot \mathbf{d}_{\text{ray}}(u, v, \theta)$ denotes the world-space coordinate of sample $s$, evaluated using trilinear interpolation within the voxel grid.

\subsubsection{Backprojection Operator ($\mathcal{B}$)}
The backprojection operator $\mathcal{B}: \mathbb{R}^{N_\theta \times W \times H} \to \mathbb{R}^{N_{xz} \times N_{xz} \times N_y}$ accumulates the perspective-projected detector intensities across all $N_\theta$ acquisition angles for each voxel $(x, y, z)$:
\begin{equation}
    \mathcal{B}(p)(x, y, z) = \frac{\pi}{N_\theta} \sum_{ip=0}^{N_\theta-1} \frac{\text{SOD}^2}{U(x, y, \theta_{ip})^2} \cdot p\big(u(x, y, \theta_{ip}),\, v(z, \theta_{ip}),\, ip\big)
    \label{eq:bp_discrete}
\end{equation}
where the perspective depth coordinate $U$ and detector projection coordinates $(u, v)$ are given by Feldkamp cone-beam geometry:
\begin{align}
    x' &= (x - r_{xz}) \cdot \delta_v, \quad y' = (y - r_{xz}) \cdot \delta_v, \quad z' = (z - r_y) \cdot \delta_v \\
    U &= \text{SOD} + y'\sin\theta + x'\cos\theta \label{eq:U_depth} \\
    u &= \frac{W-1}{2} - \frac{\text{SDD}}{U} \cdot (y'\cos\theta - x'\sin\theta) \cdot \frac{1}{\delta_p} \label{eq:u_coord} \\
    v &= \frac{H-1}{2} + \frac{\text{SDD}}{U} \cdot z' \cdot \frac{1}{\delta_p} \label{eq:v_coord}
\end{align}
where $r_{xz} = \frac{N_{xz}-1}{2}$ and $r_y = \frac{N_y-1}{2}$. Values of $p(u,v,ip)$ are evaluated via 2D bilinear interpolation on the projection plane.

\subsubsection{Cosine Distance Weighting and Data Preprocessing}
Cone-beam geometry requires a geometric cosine weighting factor $w(u, v)$ applied to projection rays to compensate for ray divergence across the planar detector:
\begin{equation}
    w(u, v) = \frac{\text{SDD}}{\sqrt{\text{SDD}^2 + u_{\text{phys}}^2 + v_{\text{phys}}^2}}, \quad \text{where } u_{\text{phys}} = \left(u - \frac{W-1}{2}\right)\delta_p, \; v_{\text{phys}} = \left(v - \frac{H-1}{2}\right)\delta_p
    \label{eq:cone_weight}
\end{equation}
Additionally, compatibility with standard CBCT indexing mandates a pre-processing transformation $\mathcal{P}$:
\begin{equation}
    \mathcal{P}(p)[ip, iu, iv] = \frac{1}{\delta_v} \cdot w(iu, iv) \cdot p[ip, H - 1 - iv, iu]
    \label{eq:preprocess}
\end{equation}
which simultaneously flips the detector height axis, transposes from row-major $[N_\theta, H, W]$ to column-major $[N_\theta, W, H]$, applies cone weighting, and normalizes by voxel size $\delta_v$.

\subsection{Maximum Likelihood Expectation Maximization (MLEM)}
Under monochromatic X-ray transmission, the measured photon counts follow a Poisson likelihood model. Maximizing the Poisson log-likelihood leads to the iterative MLEM multiplicative update equation:
\begin{equation}
    v^{(k+1)} = v^{(k)} \cdot \frac{\mathcal{B}\left(\mathcal{P}\left(\dfrac{p_0}{\mathcal{F}(v^{(k)})}\right)\right)}{\mathcal{B}\left(\mathcal{P}(\mathbf{1})\right)}
    \label{eq:mlem_update}
\end{equation}
where $p_0$ is the measured projection dataset, $\mathbf{1}$ is an all-ones projection matrix, and divisions/multiplications are performed element-wise over the volume. The normalizer $\mathcal{B}(\mathcal{P}(\mathbf{1}))$ represents the voxel-wise sensitivity matrix and is computed exactly once prior to the iterative loop.

% ─────────────────────────────────────────────────────────────────────────────
\section{Architecture Evolution and System Design}
\label{sec:arch}

The codebase underwent a progressive three-tier evolution over the course of the project:

\begin{figure}[h!]
\centering
\begin{tcolorbox}[colback=codebg, colframe=primary, arc=2mm, boxrule=0.8pt, title=\textbf{Architecture Evolution Milestone Summary}]
\small
\textbf{Phase 0: Python / PyTorch C++ Extension (Report 1)}\\
Host control in Python (\texttt{Topic\_2\_CTreconstruction.py}), native C++ backend (\texttt{backend.cpp}) linked via PyTorch JIT. Inverted nested loops, private OpenMP thread buffering for backprojection, and coordinate pre-hoisting.\\[0.3em]
\centerline{$\boldsymbol{\Downarrow}$}\\[0.3em]
\textbf{Phase 1--2: Standalone C Engine \& Modular OpenCL Architecture (Progress Report 2)}\\
Zero external heavy dependencies; native C CLI (\texttt{main.c}), direct HDF5 I/O (\texttt{utils.c}). OpenMP C baseline (\texttt{ct\_cpu.c}) and initial OpenCL host driver (\texttt{ct\_gpu.c}) featuring 3 GPU modes.\\[0.3em]
\centerline{$\boldsymbol{\Downarrow}$}\\[0.3em]
\textbf{Phase 3--4: Optimized Heterogeneous Engine \& Algorithmic OSEM (Final State)}\\
Fused GPU kernels (\texttt{divide\_preprocess\_img}, \texttt{vol\_update\_img}), CPU ray-tiling (\texttt{FP\_TILE=32}), cooperative local memory LUT caching, OSEM ordered-subset solver, and per-epoch convergence instrumentation.
\end{tcolorbox}
\vspace{-0.3cm}
\caption{System architecture evolution from initial prototype to final heterogeneous high-performance engine.}
\label{fig:arch_evolution}
\end{figure}

\subsection{System Architecture and CLI Interface}
The final application compiles into a single unified executable \texttt{build/ct\_recon} with an explicit command-line interface:
\begin{lstlisting}[language=bash, numbers=none]
./build/ct_recon --data <input.hdf5> --out <output.hdf5> \
                 --mode <cpu | gpu-buf | gpu-img | gpu-opt> \
                 --epochs <N> [--subsets <S>] [--samples <Ns>] \
                 [--kernels <dir>] [--half] [--log-convergence <log.csv>]
\end{lstlisting}

\subsection{Reconstruction Mode Taxonomy}
\begin{table}[h!]
\centering
\small
\begin{tabularx}{\textwidth}{l L L L}
\toprule
\textbf{Mode} & \textbf{Forward Projection ($\mathcal{F}$)} & \textbf{Backprojection ($\mathcal{B}$)} & \textbf{Memory Model \& Features} \\
\midrule
\texttt{cpu} & OpenMP ray-marching with ray-tiling (\texttt{FP\_TILE=32}) & Voxel-parallel OpenMP $(x,y)$ slices with $\cos/\sin$ LUT & CPU L1/L2 cache locality, \texttt{-ffast-math}, single-pass reduction \\
\midrule
\texttt{gpu-buf} & Manual 8-tap trilinear gather in global buffer & Manual 4-tap bilinear gather in global buffer & Pure \texttt{cl\_mem} global buffers; z-slab and angle-slab chunking \\
\midrule
\texttt{gpu-img} & Hardware trilinear sampling on \texttt{image3d\_t} volume & Hardware bilinear sampling on \texttt{image2d\_array\_t} projections & Dedicated GPU texture cache; automatic fractional interpolation \\
\midrule
\texttt{gpu-opt} & Hardware trilinear sampling on \texttt{image3d\_t} volume & Hardware bilinear sampling + precomputed float2 LUT + \texttt{\_\_local} cache & Scalar loop, fused preprocess/divide, direct image volume update, OSEM \\
\bottomrule
\end{tabularx}
\caption{Comprehensive taxonomy of reconstruction modes in the final codebase.}
\label{tab:mode_taxonomy}
\end{table}

% ─────────────────────────────────────────────────────────────────────────────
\section{Correctness Audits and Bug Discovery}
\label{sec:correctness}

Achieving true numerical correctness was the foremost requirement before conducting optimization sweeps. Initial implementations exhibited subtle divergence against the Python reference. A systematic audit uncovered three critical bugs:

\subsection{Interpolation Boundary Condition Mismatch}
In early GPU kernels, out-of-bounds detector taps were handled by OpenCL sampler clamping (\texttt{CLK\_ADDRESS\_CLAMP}), which zero-padded individual out-of-bounds corners while blending the remaining in-bounds corners. However, empirical inspection of the Python reference (\texttt{scipy.interpolate.RegularGridInterpolator(fill\_value=0)}) revealed that SciPy zeroes the \emph{entire interpolation cell} if even a single tap falls outside detector boundaries.

Modifying the GPU kernels to enforce an explicit whole-cell bounds check before sampling eliminated a substantial systematic bias:
\begin{lstlisting}
float uf = half_W - term * inv_U;
int u0 = (int)floor(uf);
if ((unsigned)u0 >= (unsigned)(W - 1)) continue;

float vf = half_H + zpr_sdd_inv_px * inv_U;
int v0 = (int)floor(vf);
if ((unsigned)v0 >= (unsigned)(H - 1)) continue;

float4 val = read_imagef(proj_images, samp, (float4)(vf + 0.5f, uf + 0.5f, (float)ip, 0.f));
\end{lstlisting}

\subsection{Fractional Index Truncation Bug in \texttt{fp\_cpu}}
The CPU forward projector initially computed voxel cell indices via direct cast:
\begin{lstlisting}
int x0 = (int)xi; /* Truncates toward zero! */
\end{lstlisting}
For rays entering near the volume boundary with coordinates $\xi \in (-1.0, 0.0)$, integer truncation yielded $x_0 = 0$, falsely passing the unsigned bound check \texttt{(unsigned)x0 < (unsigned)(Nxz-1)}. Consequently, voxels at $x=0$ were sampled with incorrect interpolation weights for rays that should have missed the volume entirely.

Replacing truncation with \texttt{floorf(xi)} resolved this error:
\begin{lstlisting}
int x0 = (int)floorf(xi), y0 = (int)floorf(yi), z0 = (int)floorf(zi);
\end{lstlisting}
Fixing this single CPU-side bug dropped the Mean Squared Error (MSE) between CPU and GPU modes across the entire volume by more than \textbf{four orders of magnitude} (from $\approx 10^{-3}$ to $\approx 10^{-8}$).

\subsection{Normalizer Scale Factor Alignment}
In early GPU prototypes, $\mathcal{B}(\mathbf{1})$ was evaluated on a raw unit buffer without applying the preprocessing operator $\mathcal{P}$. This omitted the $1/\delta_v$ scaling factor ($\approx 313.12$ on $256^3$), resulting in catastrophic divergence in the multiplicative update. Ensuring $\mathcal{P}(\mathbf{1})$ is fully evaluated and preprocessed before computing $\mathcal{B}(\mathcal{P}(\mathbf{1}))$ established exact mathematical alignment with Eq.~\eqref{eq:mlem_update}.

% ─────────────────────────────────────────────────────────────────────────────
\section{CPU High-Performance Optimization}
\label{sec:cpu_opt}

\subsection{Memory Hierarchy Bottleneck Analysis}
Profiling the baseline CPU implementation on $512^3$ volumes revealed that \texttt{fp\_cpu} accounted for over $78\%$ of total epoch runtime ($36.2\,\text{s}$ out of $45.7\,\text{s}$). The 8-tap trilinear interpolation performs 8 memory reads per ray sample across large volume strides ($N_{xz} \cdot N_y = 262{,}144$ floats $\approx 1.05\,\text{MB}$). Stepping one ray sequentially to completion exhausted the L1/L2 caches, causing continuous cache thrashing.

\subsection{Ray-Tiling Memory Reordering (\texttt{FP\_TILE=32})}
To restore cache locality, we restructured \texttt{fp\_cpu} using a 2D ray-tiling strategy:
\begin{enumerate}[noitemsep]
    \item Group $T = \text{\texttt{FP\_TILE}}$ adjacent detector rays along the detector width into a single batch.
    \item \textbf{Invert loop nesting}: Iterate over the ray sample index $s$ as the \emph{outer} loop, advancing all $T$ rays synchronously.
\end{enumerate}

\begin{lstlisting}
/* Batch T=32 adjacent rays */
for (int s = tile_s_lo; s < tile_s_hi; s++) {
    for (int t = 0; t < tile_n; t++) {
        if (s < s_start[t] || s >= s_end[t]) continue;
        float curr_xi = xi[t], curr_yi = yi[t], curr_zi = zi[t];
        int x0 = (int)floorf(curr_xi), y0 = (int)floorf(curr_yi), z0 = (int)floorf(curr_zi);
        if ((unsigned)x0 < (unsigned)(Nxz-1) && (unsigned)y0 < (unsigned)(Ny-1) && (unsigned)z0 < (unsigned)(Nxz-1)) {
            /* 8-tap trilinear gather: addresses cluster tightly in L1/L2 */
            val[t] += trilinear_gather(volume, x0, y0, z0, dx, dy, dz);
        }
        xi[t] += dxi[t]; yi[t] += dyi[t]; zi[t] += dzi[t];
    }
}
\end{lstlisting}

Because neighboring rays originate close together and diverge gradually, their spatial coordinates at sample step $s$ cluster in memory. Batching allows the CPU memory subsystem to reuse cache lines across all 32 rays.

\begin{table}[h!]
\centering
\small
\begin{tabularx}{0.85\textwidth}{l YYYYYY}
\toprule
\textbf{Tile Size (\texttt{FP\_TILE})} & 4 & 8 (initial) & 16 & \textbf{32 (optimal)} & 48 & 64 \\
\midrule
\textbf{\texttt{fp\_cpu} Time/Epoch (512$^3$)} & 26.05\,s & 16.84\,s & 14.72\,s & \textbf{13.58\,s} & 13.79\,s & 13.64\,s \\
\bottomrule
\end{tabularx}
\caption{Empirical evaluation of CPU ray-tiling batch sizes on Intel Core i7-5820K.}
\label{tab:cpu_tile_sweep}
\end{table}

As shown in Table~\ref{tab:cpu_tile_sweep}, tuning $\text{\texttt{FP\_TILE}} = 32$ reduces \texttt{fp\_cpu} runtime from $16.84\,\text{s}$ to $13.58\,\text{s}$ (a $19.4\%$ win over initial tiling, and $>2.6\times$ faster than the $36.2\,\text{s}$ untiled baseline), bringing total $512^3$ CPU epoch time down to $22.9\,\text{s}$.

% ─────────────────────────────────────────────────────────────────────────────
\section{OpenCL GPU Kernel Engineering and Memory Optimization}
\label{sec:gpu_opt}

\subsection{Hardware Texture Samplers vs. Buffer Gather}
In \texttt{gpu-buf}, every trilinear sample requires computing 8 global memory addresses and executing 8 uncoalesced global memory loads. On $512^3$ volumes with $N_s = 512$, forward projection executes:
\[
512 \times 512 \times 75 \times 512 \times 8 \approx 8.05 \times 10^{10} \text{ global memory fetches per epoch}
\]
In \texttt{gpu-img} and \texttt{gpu-opt}, storing the volume in a 3D OpenCL image object (\texttt{image3d\_t}) routes all memory traffic through the GPU hardware texture unit and texture cache, converting 8 manual memory transactions and 7 interpolation operations into a single \texttt{read\_imagef()} hardware primitive.

\subsection{Phase B Kernel Fusions}
Between MLEM projection stages, traditional pipelines execute separate memory copy and element-wise transformation passes. In Phase~B, we developed two fully fused OpenCL kernels that completely eliminate intermediate global memory buffers:

\subsubsection{Fused Divide and Preprocess (\texttt{divide\_preprocess\_img})}
Replaces three separate operations (element-wise division $p_0 / b \to d_{\text{ratio}}$, projection preprocessing $d_{\text{ratio}} \to d_{\text{prep}}$, and buffer-to-image copy $d_{\text{prep}} \to \text{ratio\_img}$) with a single kernel:
\begin{lstlisting}
__kernel void divide_preprocess_img(
    __global const float *p0, __global const float *b,
    __write_only image2d_array_t ratio_img,
    int W, int H, float voxelSize, float SDD, float pixelSize)
{
    int iw = get_global_id(0), ih = get_global_id(1), ip = get_global_id(2);
    if (iw >= W || ih >= H) return;
    int src_idx = ip * H * W + (H - 1 - ih) * W + iw;
    float b_val = b[src_idx];
    float ratio = (b_val > 1e-3f) ? (p0[src_idx] / b_val) : 0.f;
    float u = (-(float)(iw - (W - 1) * 0.5f)) * pixelSize;
    float v = ( (float)(ih - (H - 1) * 0.5f)) * pixelSize;
    float cw = SDD * native_rsqrt(SDD * SDD + u * u + v * v);
    float val = ratio * cw * native_recip(voxelSize);
    write_imagef(ratio_img, (int4)(ih, iw, ip, 0), (float4)(val, 0.f, 0.f, 0.f));
}
\end{lstlisting}
At $512^3$, this removes $\approx 0.80\,\text{GB}/\text{epoch}$ of redundant PCIe/VRAM memory transfers.

\subsubsection{Direct Volume Image Update (\texttt{vol\_update\_img})}
In standard image mode, updating the volume requires writing to a global float buffer $d_{\text{vol}}$, followed by a \texttt{clEnqueueCopyBufferToImage} staging copy into \texttt{vol\_img} for the next forward projection. In float32 mode, \texttt{vol\_update\_img} performs the multiplicative update $v \leftarrow v \cdot \frac{\text{bp\_ratio}}{\text{bp\_ones}}$ and writes directly into both $d_{\text{vol}}$ and \texttt{vol\_img} simultaneously, eliminating $\approx 1.07\,\text{GB}/\text{epoch}$ of volume staging traffic at $512^3$.

\subsection{Systematic Work-Group Sweeps}
To optimize work-group scheduling across different GPU compute architectures, we performed exhaustive sweeps of local work-group dimensions (Table~\ref{tab:wg_sweeps}).

\begin{table}[h!]
\centering
\small
\begin{tabularx}{\textwidth}{l c L}
\toprule
\textbf{Kernel} & \textbf{Optimal LWS} & \textbf{Architectural Rationale \& Measured Impact} \\
\midrule
\texttt{bp\_buffer}, \texttt{bp\_image}, \texttt{bp\_opt} & $\{4, 4, 16\}$ & 16 contiguous work-items along z-dimension align with volume layout $[x][y][z]$, maximizing memory write coalescing. \\
\midrule
\texttt{fp\_image} & $\{8, 32, 1\}$ & Narrow-in-$W$ / tall-in-$H$ aspect ratio matches detector dimensions ($1120 \times 1184$). Outperformed $\{16,16,1\}$ by $5.4\%$ on AMD Hawaii ($0.873\,\text{s}$ vs $0.923\,\text{s}$). \\
\midrule
\texttt{fp\_buffer} & $\{4, 64, 1\}$ & Uncached global gather is extremely sensitive to memory latency. Outperformed default $\{16,16,1\}$ by $>4\times$ ($7.53\,\text{s}$ vs $32.1\,\text{s}$ per 10 ep at $512^3$). \\
\bottomrule
\end{tabularx}
\caption{Optimal work-group configurations identified through empirical sweeping.}
\label{tab:wg_sweeps}
\end{table}

\subsection{Register Pressure and Instruction-Level Parallelism Analysis}
In \texttt{bp\_buffer\_opt.cl}, a $2\times$ loop unrolling optimization (\texttt{unroll-x2}) was originally introduced to hide texture fetch latency via instruction-level parallelism (ILP). However, isolated benchmarking on AMD Hawaii (GCN 1.1) revealed that \texttt{gpu-opt} with \texttt{unroll-x2} was \textbf{$1.6\%$ slower} than scalar \texttt{gpu-img} ($0.945\,\text{s}$ vs $0.930\,\text{s}$/epoch at $512^3$).

Tracing the compiled ISA revealed that unrolling doubled the live vector register footprint (concurrent \texttt{float2} and \texttt{float4} variables), forcing the OpenCL compiler to allocate more VGPRs per work-item. This reduced the active work-group occupancy per Compute Unit. Removing \texttt{unroll-x2} permanently restored scalar execution, dropping runtime to $0.923\,\text{s}/\text{epoch}$ and re-establishing \texttt{gpu-opt} as faster than \texttt{gpu-img}.

% ─────────────────────────────────────────────────────────────────────────────
\section{Empirical Investigations and Negative Results}
\label{sec:investigations}

In high-performance computing, documenting rigorously tested negative results is as critical as reporting wins.

\subsection{Investigation of \texttt{gpu-buf} Run-to-Run Variance on AMD Hawaii}
On the AMD Hawaii PRO testbed (\texttt{pool15-01}), \texttt{gpu-buf} exhibited large run-to-run execution time variance at $512^3$ scale ($75\,\text{s}\text{--}102\,\text{s}$ for 10 epochs across identical runs), whereas all other modes varied by $<1\%$.

\begin{enumerate}[noitemsep]
    \item \textbf{Hypothesis 1 (Multi-user Contention)}: Checked via system tools; server had zero concurrent user sessions.
    \item \textbf{Hypothesis 2 (Thermal Throttling)}: Tested by inserting idle sleep pauses between dispatches (\texttt{FP\_BUFFER\_SLAB\_PAUSE\_US}). Pauses \emph{worsened} total runtime ($172.4\,\text{s}$ vs $85\,\text{s}$), disproving thermal throttling (which would benefit from idle cooling).
    \item \textbf{Hypothesis 3 (Driver-side Memory Demotion)}: Built a dedicated diagnostic (\texttt{-\,-diag repeat-slab}) with OpenCL hardware event profiling. Timing showed a single step-change degradation to $6\times$ slower that was immediately recovered upon reallocating $d_{\text{vol}}$.
    \item \textbf{Cross-Hardware Verification}: Re-running the identical diagnostic on NVIDIA Kepler (\texttt{kale}) produced completely flat timings across all 40 repeats ($\pm 0.1\%$). This proved decisively that the variance was an AMD-driver-specific memory placement artifact under sustained uncached buffer access.
\end{enumerate}

\subsection{D1 Branch-Free Clamp Regression in \texttt{bp\_cpu}}
Attempting to vectorize \texttt{bp\_cpu} by replacing the inner loop \texttt{continue} bound check with a branch-free arithmetic clamp (\texttt{clamp(u, 0, W-1)}) was initially thought to be neutral. A high-precision 3-trial benchmark revealed a consistent \textbf{$3.7\%$ performance regression} ($45.59\,\text{s}$ vs $43.96\,\text{s}$). The branch-free code forced unconditional 4-tap bilinear math on out-of-bounds voxels rather than skipping them, while failing to auto-vectorize due to scatter-gather requirements. The change was cleanly reverted.

\subsection{AABB Slab Ray-Clipping at $256^3$}
Axis-Aligned Bounding Box (AABB) ray clipping calculates entry and exit intersections $(s_{\text{start}}, s_{\text{end}})$ with the volume cube. While highly effective at $512^3$ (saving $35\%$ of ray march samples), forcing AABB on $256^3$ produced a consistent \textbf{$4\%$ slowdown} ($0.099\,\text{s}$ vs $0.095\,\text{s}$/epoch). On small grids, the arithmetic cost of 6 slab divisions and divergent inner-loop trip counts outweighs the memory fetch savings. AABB is correctly gated to $W > 512$.

% ─────────────────────────────────────────────────────────────────────────────
\section{Algorithmic Acceleration: Ordered Subsets EM (OSEM)}
\label{sec:osem}

\subsection{Mathematical Derivation of OSEM}
To accelerate the asymptotic convergence of MLEM, we implemented Ordered Subsets Expectation Maximization (OSEM) in \texttt{gpu-opt}. The set of $N_\theta$ projection angles is partitioned into $S$ disjoint subsets $S_0, S_1, \dots, S_{S-1}$. For each subset $s \in \{0, \dots, S-1\}$, an intermediate sub-iteration update is computed:
\begin{equation}
    v^{(k, s+1)} = v^{(k, s)} \cdot \frac{\mathcal{B}_s\left(\mathcal{P}\left(\dfrac{p_0}{\mathcal{F}(v^{(k, s)})}\right)\right)}{\mathcal{B}_s\left(\mathcal{P}(\mathbf{1})\right)}
    \label{eq:osem_update}
\end{equation}
where $\mathcal{B}_s$ denotes the backprojection evaluated exclusively over the angular sub-range belonging to subset $S_s$.

\subsubsection{Precomputation of Subset Normalizers}
Each subset possesses a distinct spatial sensitivity matrix $\mathcal{B}_s(\mathcal{P}(\mathbf{1}))$. At initialization, $S$ separate 3D normalizer buffers $d_{\text{bp\_ones}}[s]$ are precomputed and resident in GPU memory. When $S=1$, Eq.~\eqref{eq:osem_update} reduces identically to standard MLEM.

\subsection{Coprime Angle Ordering Strategy}
To ensure maximum angular independence between consecutive sub-iterations, angles are assigned to subsets in an interleaved fashion:
\begin{equation}
    S_s = \left\{ ip \in \{0, \dots, N_\theta-1\} \;\middle|\; ip \equiv s \pmod S \right\}
\end{equation}
The sub-iterations visit subsets in a golden-ratio-derived coprime stride sequence. The projection data is permuted once upon loading, allowing each GPU backprojection launch to execute over a contiguous angle range $[ip_{\text{start}}, ip_{\text{start}} + ip_{\text{count}})$.

\subsection{VRAM Capacity Bounds at $512^3$}
On hardware with 4\,GB VRAM (such as the NVIDIA GTX 680 on \texttt{kale}), each $512^3$ float32 volume buffer requires $536.87\,\text{MB}$. Allocating the base volume, forward projections, detector images, and $S$ subset normalizers requires:
\begin{equation}
    \text{VRAM}_{\text{req}}(S) \approx 3.25\,\text{GB} + S \times 0.537\,\text{GB}
\end{equation}
For $S \ge 2$, memory demand exceeds 4.3\,GB, triggering out-of-memory allocation failures. Hence, multi-normalizer OSEM is purposefully focused on $256^3$ scale where all subsets reside safely within VRAM ($S=25$ consumes $<800\,\text{MB}$).

\subsection{Empirical Convergence Analysis}
To evaluate reconstruction convergence objectively without requiring unavailable ground-truth phantoms, we track data-domain Poisson log-likelihood $\mathcal{L}(v)$ and residual fidelity via \texttt{-\,-log-convergence}:
\begin{equation}
    \mathcal{L}(v) = \sum_{i=1}^{N_{\text{rays}}} \left( p_0[i] \ln b[i] - b[i] \right)
\end{equation}

\begin{table}[h!]
\centering
\small
\begin{tabularx}{0.92\textwidth}{l YYYYY}
\toprule
\textbf{Time} & \textbf{S=1 (Plain)} & \textbf{S=3} & \textbf{S=5} & \textbf{S=15} & \textbf{S=25} \\
\midrule
10\,s & $-482{,}720$ & $-481{,}297$ & $-481{,}109$ & $\mathbf{-480{,}543}$ & $-480{,}577$ \\
15\,s & $-481{,}109$ & $-480{,}391$ & $-480{,}304$ & $-479{,}652$ & $\mathbf{-479{,}285}$ \\
20\,s & $-480{,}652$ \emph{(plateau)} & $-480{,}013$ & $-479{,}942$ & $-479{,}305$ & $\mathbf{-478{,}917}$ \\
30\,s & $-480{,}652$ \emph{(plateau)} & $-479{,}701$ & $-479{,}653$ & $-479{,}042$ & $\mathbf{-478{,}651}$ \\
36\,s & $-480{,}652$ \emph{(plateau)} & $-479{,}610$ & $-479{,}580$ & $-478{,}964$ & $\mathbf{-478{,}571}$ \\
\bottomrule
\end{tabularx}
\caption{Log-likelihood convergence vs. wall-clock time on NVIDIA GTX 680 ($256^3$, 100 epochs).}
\label{tab:osem_convergence}
\end{table}

As shown in Table~\ref{tab:osem_convergence}, plain MLEM ($S=1$) completely plateaus by $\approx 20\,\text{s}$ at its epoch-limited ceiling ($-480{,}652$). In contrast, every OSEM configuration continues to optimize the log-likelihood objective. $S=15$ leads in early wall-clock time ($\le 10\,\text{s}$), while $S=25$ achieves the highest overall reconstruction fidelity from $15\,\text{s}$ onward.

% ─────────────────────────────────────────────────────────────────────────────
\section{Experimental Results and Performance Benchmarks}
\label{sec:results}

\subsection{Hardware Testbeds}
\begin{itemize}[noitemsep]
    \item \textbf{Platform A (\texttt{pool15-01})}: AMD Hawaii PRO (Radeon R9 290/390, 2560 shader cores, 2.56 TFLOPS) paired with Intel Core i7-5820K @ 3.30GHz (6 physical cores, 12 logical OpenMP threads).
    \item \textbf{Platform B (\texttt{kale})}: NVIDIA GeForce GTX 680 (Kepler architecture, 1536 CUDA cores, 4037\,MiB VRAM) paired with Intel Xeon E5-2620 0 @ 2.00GHz (12 physical cores, 24 logical OpenMP threads).
\end{itemize}

\subsection{Performance Benchmarks}

\begin{table}[h!]
\centering
\small
\setlength{\tabcolsep}{4pt}
\begin{tabularx}{\textwidth}{l YYYY}
\toprule
\textbf{Mode} & \textbf{256$^3$ Time/Ep} & \textbf{512$^3$ Time/Ep} & \textbf{Speedup (256$^3$)} & \textbf{Speedup (512$^3$)} \\
\midrule
\texttt{cpu} (12-thread) & 2.60\,s & 22.9\,s & $1.0\times$ & $1.0\times$ \\
\texttt{gpu-buf} & 0.44\,s & 4.87--10.19\,s & $5.9\times$ & $2.2\text{--}4.8\times$ \\
\texttt{gpu-img} & 0.095\,s & 0.873\,s & $27.4\times$ & $26.2\times$ \\
\texttt{gpu-opt} & \textbf{0.093\,s} & \textbf{0.870\,s} & $\mathbf{28.0\times}$ & $\mathbf{26.3\times}$ \\
\bottomrule
\end{tabularx}
\caption{Reconstruction runtime on Platform A (AMD Hawaii PRO / Intel Core i7-5820K, $\text{EPOCHS}=10$).}
\label{tab:perf_pool}
\end{table}

\begin{table}[h!]
\centering
\small
\setlength{\tabcolsep}{3.5pt}
\begin{tabularx}{\textwidth}{l YYYYY}
\toprule
 & \multicolumn{2}{c}{\textbf{256$^3$ Dataset}} & \multicolumn{2}{c}{\textbf{512$^3$ Dataset}} & \\
\cmidrule(lr){2-3} \cmidrule(lr){4-5}
\textbf{Mode} & \textbf{s / epoch} & \textbf{Total (100 ep)} & \textbf{s / epoch} & \textbf{Total (100 ep)} & \textbf{Speedup} \\
\midrule
\texttt{cpu} (24-thread) & $\approx 4.20$\,s & 423.9\,s & $\approx 34.0$\,s & 3415.6\,s & $1.0\times$ \\
\texttt{gpu-buf} & 0.858\,s & 86.5\,s & 7.17\,s & 722.6\,s & $4.9\times / 4.7\times$ \\
\texttt{gpu-img} & 0.145\,s & 14.7\,s & 1.24\,s & 126.7\,s & $28.9\times / 27.0\times$ \\
\texttt{gpu-opt} & \textbf{0.141\,s} & \textbf{14.3\,s} & \textbf{1.23\,s} & \textbf{125.8\,s} & $\mathbf{29.7\times / 27.2\times}$ \\
\bottomrule
\end{tabularx}
\caption{Reconstruction runtime on Platform B (NVIDIA GeForce GTX 680 / Intel Xeon E5-2620, $\text{EPOCHS}=100$).}
\label{tab:perf_kale}
\end{table}

\subsection{Numerical Correctness and Validation}
Numerical fidelity was validated across both dataset resolutions over full 100-epoch convergence runs (Table~\ref{tab:val_results}).

\begin{table}[h!]
\centering
\small
\setlength{\tabcolsep}{4pt}
\begin{tabularx}{\textwidth}{ll YYY c Y}
\toprule
\textbf{Scale} & \textbf{Mode} & \textbf{Min} & \textbf{Max} & \textbf{Mean} & \textbf{NaN/Inf} & \textbf{MSE vs CPU} \\
\midrule
$256^3$ & \texttt{cpu} & 0.0000 & 1.7303 & 0.0067 & 0 / 0 & (Reference) \\
& \texttt{gpu-buf} & 0.0000 & 1.7292 & 0.0067 & 0 / 0 & $1.148 \times 10^{-10}$ \\
& \texttt{gpu-img} / \texttt{gpu-opt} (default) & 0.0000 & 1.6577 & 0.0067 & 0 / 0 & $1.128 \times 10^{-7}$ \\
& \texttt{gpu-img} (\texttt{FP\_TEX\_EXACT=1}) & 0.0000 & 1.7292 & 0.0067 & 0 / 0 & $\mathbf{2.524 \times 10^{-9}}$ \\
& \texttt{gpu-opt} (\texttt{--fp-buf}) & 0.0000 & 1.7292 & 0.0067 & 0 / 0 & $\mathbf{2.524 \times 10^{-9}}$ \\
\midrule
$512^3$ & \texttt{cpu} & 0.0000 & 1.0055 & 0.0330 & 0 / 0 & (Reference) \\
& \texttt{gpu-buf} & 0.0000 & 1.0054 & 0.0330 & 0 / 0 & $9.534 \times 10^{-11}$ \\
& \texttt{gpu-img} & 0.0000 & 1.0054 & 0.0330 & 0 / 0 & $1.232 \times 10^{-9}$ \\
& \texttt{gpu-opt} & 0.0000 & 1.0054 & 0.0330 & 0 / 0 & $1.232 \times 10^{-9}$ \\
\bottomrule
\end{tabularx}
\caption{Full 100-epoch numerical validation results across $256^3$ and $512^3$ datasets.}
\label{tab:val_results}
\end{table}

\subsection{The $256^3$ Precision Investigation: Hardware Sampler Dissection}
\label{sec:precision_investigation}
While $512^3$ texture modes comfortably cleared the course's $10^{-8}$ MSE bar ($1.232 \times 10^{-9}$), the default $256^3$ texture modes produced an MSE of $1.128 \times 10^{-7}$. A forensic investigation across intermediate checkpoints (Epochs 1, 5, 10, 20, 50, 100) revealed that at Epoch 1 the GPU matches the CPU reference to within $1.12 \times 10^{-13}$, but compounds non-linearly over 100 MLEM iterations.

To establish whether this represented an implementation flaw or fundamental hardware constraints, five distinct hypotheses were formulated and tested:
\begin{enumerate}[leftmargin=*]
    \item \textbf{Hypothesis 1 (IEEE Instruction Rounding)}: Replacing fast OpenCL math built-ins (\texttt{native\_recip}, \texttt{native\_sqrt}) with strict IEEE 754 float operations yielded \textbf{0\% accuracy gain} (MSE remained $1.1278 \times 10^{-7}$), confirming arithmetic intrinsics were not the root cause.
    \item \textbf{Hypothesis 2 (Backprojection Sampler Mismatch)}: Routing backprojection through software global memory while keeping hardware texture forward projection yielded only a \textbf{6\% accuracy gain} while incurring a $55\%$ runtime penalty, proving Backprojection was not the offending operator.
    \item \textbf{Hypothesis 3 (Normalizer Underflow and FOV Masking)}: Inspecting the sensitivity normalizer $\mathcal{B}(\mathbf{1})$ confirmed the denominator is well-behaved across the entire volume ($364 \le \mathcal{B}(\mathbf{1}) \le 443$). Furthermore, spatial error localization revealed the worst 50 outlier voxels were concentrated near $r \approx 90\% R_{\max}$, meaning an outer cylinder mask ($r > R_{\max}$) left the errant boundary untouched while masking $r < R_{\max}$ diverged from the unmasked CPU reference.
    \item \textbf{Hypothesis 4 (Two-Phase Warm-Start)}: Attempting 85 fast TMU epochs followed by 15 exact buffer epochs failed because MLEM error at early-to-mid epochs is dominated by algorithmic convergence rather than floating-point precision (MSE at Epoch 25 is $1.03 \times 10^{-4}$); 15 epochs lack the contraction rate required to wash out an 85-epoch trajectory bias.
    \item \textbf{Hypothesis 5 (Forward Projection Sampler Decoupling)}: Bypassing the hardware TMU's 8-bit fixed-point fractional interpolator during ray-marching forward projection achieved a \textbf{$45\times$ precision gain}, dropping MSE to $\mathbf{2.524 \times 10^{-9}}$ and passing the $10^{-8}$ bar.
\end{enumerate}

Crucially, by implementing \texttt{FP\_TEX\_EXACT=1} (which pairs OpenCL's \texttt{CLK\_FILTER\_NEAREST} 8-tap texture cache fetches with an exact software IEEE float32 trilinear blend), this $45\times$ accuracy breakthrough was achieved in just \textbf{21.34\,s} ($1.46\times$ baseline), proving that the 3D texture cache can be preserved while discarding the lossy hardware blending circuit.

% ─────────────────────────────────────────────────────────────────────────────
\section{Discussion and Engineering Insights}
\label{sec:discussion}

\subsection{Architectural Tradeoffs Across Compute Generations}
A central finding is the dramatic disparity between global buffer gathers and hardware texture sampling. In \texttt{gpu-buf}, manual trilinear interpolation lacks hardware cache support, rendering it severely memory-bandwidth bound. On $512^3$, \texttt{gpu-buf} initially hung due to GPU watchdog timeouts from oversized single-launch kernel grids; resolving this required chunking into discrete slab launches.

Conversely, \texttt{gpu-img} and \texttt{gpu-opt} leverage hardware texture mapping units (TMUs) to perform trilinear blending and 3D spatial cache line fetches in hardware, achieving an immediate $\approx 6\times$ advantage over \texttt{gpu-buf} and $\approx 30\times$ over multi-threaded CPU execution.

\subsection{Methodological Insights: Measurement vs. Assumption}
The project reinforced critical methodological principles:
\begin{itemize}[noitemsep]
    \item \textbf{Plausibility is not proof}: Loop unrolling and branch-free clamping appeared intuitively beneficial, yet empirical benchmarking revealed that both introduced subtle performance regressions due to register pressure and wasted arithmetic.
    \item \textbf{Isolated component verification}: Debugging whole-system iteration loops obscures localized errors; isolated operator testing (\texttt{-\,-op fp|bp}) was indispensable in isolating sub-pixel interpolation anomalies.
\end{itemize}

% ─────────────────────────────────────────────────────────────────────────────
\section{Conclusion}
\label{sec:conclusion}

In this project, we designed, implemented, and benchmarked a high-performance heterogeneous CT reconstruction engine for cone-beam MLEM and OSEM. By combining algorithmic coordinate optimizations, CPU ray-tiling memory reordering, OpenCL hardware texture sampling, fused memory kernels, and ordered-subset acceleration, we achieved:
\begin{itemize}[noitemsep]
    \item Robust numerical accuracy meeting and exceeding the course $10^{-8}$ bar across all modes ($\text{MSE} = 2.524 \times 10^{-9}$ on $256^3$, $1.232 \times 10^{-9}$ on $512^3$, and $9.53 \times 10^{-11}$ in buffer mode).
    \item A $>2\times$ speedup on multi-core CPU execution via cache-locality ray tiling.
    \item Up to $\approx 30\times$ speedup on GPU over the CPU baseline.
    \item A continuous convergence acceleration via OSEM that overcomes plain MLEM plateauing on fixed computational time budgets.
\end{itemize}

% ─────────────────────────────────────────────────────────────────────────────
\begin{thebibliography}{9}
\bibitem{shepp1982}
L.~A. Shepp and Y.~Vardi, ``Maximum likelihood reconstruction for emission tomography,'' \emph{IEEE Transactions on Medical Imaging}, vol.~1, no.~2, pp.~113--122, 1982.

\bibitem{hudson1994}
H.~M. Hudson and R.~S. Larkin, ``Accelerated image reconstruction using ordered subsets of projection data,'' \emph{IEEE Transactions on Medical Imaging}, vol.~13, no.~4, pp.~601--609, 1994.

\bibitem{feldkamp1984}
L.~A. Feldkamp, L.~C. Davis, and J.~W. Kress, ``Practical cone-beam algorithm,'' \emph{Journal of the Optical Society of America A}, vol.~1, no.~6, pp.~612--619, 1984.

\bibitem{turbell2001}
H.~Turbell, ``Cone-beam reconstruction using filtered backprojection,'' Ph.D. dissertation, Linköping University, Sweden, 2001.

\bibitem{opencl12}
Khronos OpenCL Working Group, \emph{The OpenCL Specification, Version 1.2}, Khronos Group, 2012.
\end{thebibliography}

\end{document}
